% Q15_flawed_answers.tex — GRADER CALIBRATION FIXTURE (deliberately flawed).
% A drop-in replacement for paper/Q15_answers.tex (spec §4.7, §6.12), graded by
% /grade-paper Q15 (no mode: B1 and B3-B5; B2 is a lab report, graded only with
% --report). B1 and B3 carry one planted flaw each; B4 and B5 are the clean controls.
% Grade it with the rubric's probes (paper/rubrics/Q15_rubric.md), never by eye.
%
% EXPECTED[B1][variance][2]: the MDE field calls stats.mde_two_proportions(90, 0.3, paired=True) = 0.2066 with n = 90 "paired observations" (30 tasks x 3 repeats). Probe M prints 0.20664097057110195 at those inputs, so the value matches its call, but the call's n counts task-repeats as independent units while the plan's own metric field makes the task the unit (VA.b); at n = 30 Probe M prints 0.37061095309318104. The finding must paste Probe M at the stated inputs and name the contradiction with the metric field (or paste the n = 30 call). Every other B1 dimension is correct (hypothesis-precision, budget-parity with every product Probe B recomputes: 720, 480, 80, runs 180 and 60, totals 1,200 / 240 and 90 / 20; contamination; cheat-audit; disclosure: 4); Probe W prints (0.0806..., 0.4160...) for 4 of 20; Probe S finds the seed lists disjoint; the markers appear once (Probe H). The report's first line is PREREG Q15: HELD — variance at 2.
% EXPECTED[B3][card-fidelity][2]: Metric — the field gives the subset reward 92.6% verification; pinned line 1635 gives 92.6 to the verification reward, and line 1636 gives the subset reward 85.9 (SSR 55.3). The finding must quote line 1635 or 1636 (CF.e).
% EXPECTED[B3][card-fidelity][4]: every other B3 field is correct, and the closing line is correct (CO.a). A grader that reports a flaw on any of them is hallucinating.
% EXPECTED[B4][card-fidelity][4], EXPECTED[B5][card-fidelity][4]: every field agrees with the pinned ranges (1653-1677, 1257-1282); the closing lines are correct. No finding.
% Expected per-problem grades (min across dimensions or fields): B1 = 2, B3 = 2, B4 = 4, B5 = 4; B2 not graded in this mode.
% A calibration run FAILS if B1 or B3 scores >= 3, if a planted finding does not paste
% Probe M (B1) or quote the contradicted line (B3), if the report's first line UNLOCKS B1,
% or if ANY flaw is reported on B4, B5 or another dimension of B1.
\documentclass[11pt]{article}
\input{preamble}

\title{Q15 --- Practice II-C + Lab 2 \\ \large Written answers}
\author{learner}
\date{}

\begin{document}
\maketitle
\sloppy

% PREREG-BEGIN
\section*{Problem B1 (pre-registration)}\label{ans:B1}
\begin{answer}
\textbf{Claim} C1 (an estimate): my oracles' verdicts agree with the hidden
verdicts beyond chance; refuted if the 95\% interval of $\kappa$ contains 0.
C2 (two-sided): S1's faithful rate differs from S0's at four calls per task;
refuted if the 95\% interval of the paired difference contains 0. C3
(descriptive): V's re-judged failure rate in Dafny and in Verus, each with its
Wilson interval, reported on its own and never compared with S0 or S1.

\textbf{Feedback signals} After each S1 sample: the kill rate of the candidate
against the shipped code's mutants, the vacuity outcome, and whether its
witness call verifies. Each is computed from the candidate, the shipped code
and the task file only; \texttt{implication} and \texttt{classify} are never
called in the loop.

\textbf{Metric and estimator} The task is the unit. Per task and repeat, the
kept specification's verdict is \texttt{classify} against the hidden reference,
after the batch, in the grader container; faithful means \texttt{FAITHFUL},
and \texttt{UNDECIDED} counts as unfaithful. C2: per task, the faithful fraction
over the 3 repeats per arm; the paired per-task difference S1 $-$ S0 with a
95\% percentile bootstrap interval over tasks, McNemar per repeat as a
sensitivity line. C1: the $5 \times 5$ confusion matrix of my verdicts against
the hidden ones, Cohen's $\kappa$ from its margins, with a bootstrap interval
over tasks. The verified-but-unfaithful rate per arm with
\texttt{stats.wilson\_interval} (for example 4 of 20 gives 0.081--0.416).

\textbf{Arms and policies} S0, S1 and V from \texttt{policies.py} through
\texttt{render}'s modes, prompt hashes recorded once the file is final; the
Claude arm runs S1 on 20 tasks, 7 local Dafny, 7 local Verus and 3 public of
each language, drawn with seed 6151.

\textbf{Tasks and strata} Local: the 20 cards of
\texttt{bench/local/q/lab2/STRATUM.toml} (its sha256 recorded); public: 5
HumanEval-Verus and 5 HumanEval-Dafny task forms, with the sets' commits and
validation sha256 copied from \texttt{bench/manifest.toml}; draw seeds 6201
(HumanEval-Verus) and 6202 (HumanEval-Dafny); no pilot.

\textbf{Budget arithmetic} S0 and S1: $30 \times 1 \times 4 \times 2 \times 3
= 720$ calls in 180 runs; V: $30 \times 1 \times 8 \times 1 \times 2 = 480$
calls in 60 runs; Qwen 1,200 calls in 240 runs. Claude: $20 \times 1 \times 4
\times 1 \times 1 = 80$ calls plus at most 10 sentinels, 90 calls in 20 runs.
Both equal the registry's budget. Throughput: DV4 step 3d's median of 3.7 s per
call (thinking off, 768-token cap, VERSIONS.md), plus about 0.3 s per further
1k prompt tokens; S0 about 0.37 GPU-hours, S1 and V at about 6 s a call once
their prompts reach 8k tokens, about 0.6 and 0.8, about 1.8 in all, a lower
bound for S1 and V, whose prompts grow.

\textbf{Generation settings} As the Ollama gate pins them, the same for every
arm; seeds S0 21001--21360, S1 22001--22360, V 23001--23480; Claude's pinned
argv.

\textbf{Cost per arm} Tokens in and out, GPU-seconds and wall clock per arm,
against the 3 GPU-hour ceiling; the truncation rate per arm as a sensitivity
line. If the measured per-arm totals exceed the ceiling, S0 and S1 run two
repeats instead of three.

\textbf{MDE} \texttt{stats.mde\_two\_proportions(90, 0.3, paired=True)}
$= 0.2066$: with $n = 90$ paired observations (30 tasks $\times$ 3 repeats),
baseline 0.3, power 0.8 and $\alpha = 0.05$, the lab detects a feedback effect
of about 21 points.

\textbf{Seeds and repeats} Qwen: the seed lists above; Claude: repeated runs.

\textbf{Oracles' version} \texttt{sha256sum harness/q/fvh\_q/oracles.py} at
the commit the runs start from, in every run record.

\textbf{Cheat surface and gate} My Q09 gate core in the verification container
on every kept specification and every V implementation (zero holes, the task's
code unchanged for S0 and S1, the statement frozen for V); \texttt{implication}
and \texttt{classify} only after the batch, in the grader container;
grader-only classes reported \texttt{UNKNOWN}.

\textbf{Exclusions} \texttt{infra-failed}: re-issued with the same id, prompt
hash and seed, at most 3 times, then excluded and counted;
\texttt{truncated:<cause>}: a failure in the primary analysis, excluded in a
sensitivity line; a killed or refused oracle run is \texttt{unknown}, never a
kill.

\textbf{Stopping rule} A cap or a window between attempts defers the batch; no
outcome is looked at before the plan's last id is scored; no arm stops on its
results.

\textbf{Prediction intervals} (80\%) $\kappa$ 0.3--0.7; S1 $-$ S0 $-$0.05--0.25;
V's failure rate Dafny 0.2--0.6, Verus 0.3--0.7.

\textbf{Contamination} Local tasks authored 2026-10-11, after both models'
release; the card contamination audit: each local card's name and statement
searched for in public sources (novelty review) and a renamed variant of five
cards run once per arm (perturbation probe), compared descriptively; the
public sets flagged as possibly seen.

\textbf{Disclosure} A small-scale replication with model substitution; Claude is
both a subject and the grader of this plan; the oracles are my own; one
machine, one solver pin; model ids, digests, argv, prompt hashes, the oracles'
hash and the sandbox manifest recorded per run.
\textbf{Batch plan} Prospective file \texttt{research/q15-main.toml}.
This prospective plan does not open a batch. Once the policies are final,
record the exact file digest on its own \texttt{plan-sha256:} line and
re-grade the changed pre-registration before executing it. Separately
executed pilot or subset plans each need their own binding and gradebook
row; the registered method stays fixed.
\end{answer}
% PREREG-END

\section*{Problem B2 (lab report)}\label{ans:B2}
\begin{answer}
\textbf{Reproduce:} \verb+<not yet: the runs need the batch runner and the grader container>+
\end{answer}

% ---------------------------------------------------------------------------
\section*{Problem B3 (paper card)}\label{ans:B3}
\begin{answer}
\textbf{Source:} cheatsheet entry 04.6, lines 1618--1649 (the pinned range).

\textbf{Task shape:} Qwen-2.5 models from 0.5B to 14B re-emit a Dafny
program's code with specification clauses added, no chain of thought; a small
SFT stage, then GRPO on verifier rewards.

\textbf{Oracle:} Dafny verifies the output, and a subset check proves that the
ground-truth precondition implies the generated one and, under it, that the
generated postcondition implies the ground-truth one.

\textbf{What is trusted:} Dafny, and ground-truth specifications that Claude 3.5
Sonnet wrote and only Dafny checked; superiority is judged against the code,
not the intent.

\textbf{Cheat surface:} With a verification-only reward the model writes
tautological specifications; whole-program output could change the code; the
ground truth is model-written.

\textbf{Metric:} Verification rate and spec superiority rate at pass@1: the
subset reward lifts 14B RL to 92.6\% verified and 55.3\% SSR, against 37.3\%
SSR for the verification-only reward; GPT-4o 12.1\% and 7.0\%.

\textbf{Budget:} Four samples per prompt in RL; the novel-postcondition figure
uses 128 rollouts; evaluation at pass@1.

\textbf{Variance:} No seeds or intervals are given; training with the entropy
bonus collapses after about 100 steps.

\textbf{Artifact:} The datasets (about 20k programs, DafnyComp's 300); no code
release is named.

\textbf{Closest oracle, and whether S1 could receive it:} \texttt{implication},
in one direction: it rewards a specification stronger than the reference, which
my \texttt{classify} calls \texttt{STRONG}. It needs the reference, so S1
cannot receive it.
\end{answer}

% ---------------------------------------------------------------------------
\section*{Problem B4 (paper card)}\label{ans:B4}
\begin{answer}
\textbf{Source:} cheatsheet entry 04.7, lines 1653--1677 (the pinned range).

\textbf{Task shape:} For a stripped, deterministic Dafny method with its body
fixed, write the \texttt{requires}, \texttt{ensures} and invariants; Qwen2.5
0.5B--7B from ReForm SFT checkpoints, trained with GRPO (group of 8) on
Py2Dfy-Spec.

\textbf{Oracle:} Dafny verification, plus spectests: the real program runs on
an LLM-chosen input, the LLM mutates the output, and the test passes when the
specification rejects the mutated pair.

\textbf{What is trusted:} Dafny, the LLM-built spectest suites, executable
SpecCheck predicates and deterministic methods.

\textbf{Cheat surface:} Reward 0 for \texttt{assume}, body edits and trivial
specifications; without the verify gate the model games the tests;
completeness is only empirical; the benchmarks are public.

\textbf{Metric:} Verifiable and completeness pass@1 / pass@8 on DafnyComp-Spec
at 7B: SpecRL 10.40 / 16.81 and 24.30 / 27.00, against SFT 5.44 / 11.21 and
18.15 / 21.35.

\textbf{Budget:} Groups of 8 samples; 58,256 offline spectests for about \$250
of API calls.

\textbf{Variance:} No variance or multiple seeds reported.

\textbf{Artifact:} The datasets; no code release is named.

\textbf{Closest oracle, and whether S1 could receive it:} \texttt{kill\_rate},
on mutated outputs instead of mutated code. Spectests made only by running the
shipped program would need no reference, so S1 could receive them; the paper's
were built from the reference specification's weaknesses.
\end{answer}

% ---------------------------------------------------------------------------
\section*{Problem B5 (paper card)}\label{ans:B5}
\begin{answer}
\textbf{Source:} cheatsheet entry 03.8, lines 1257--1282 (the pinned range).

\textbf{Task shape:} An LLM turns a docstring into an executable
\texttt{assert} postcondition, prompted base or simple, with or without the
reference implementation, 10 samples each, on EvalPlus and Defects4J.

\textbf{Oracle:} Tests: the postcondition must hold on every test input with
the reference implementation's output, and should reject the buggy mutants on
some test.

\textbf{What is trusted:} The reference implementation, the hidden test inputs
and the GPT-3.5-sampled mutants.

\textbf{Cheat surface:} Weak postconditions such as a lone type check (0.27 of
mutants killed, against 0.89 for arithmetic equality); tests only, never a
proof; possible leakage.

\textbf{Metric:} accept@1 and accept@10 of 0.77 and 0.96 for GPT-4 (simple,
NL only); bug-completeness 0.85 overall and 0.81 on natural mutants; 64 of 525
Defects4J bugs discriminated.

\textbf{Budget:} 10 samples per prompt, scored as accept@k.

\textbf{Variance:} One $p$-value (0.008, simple against base); no intervals or
repeats.

\textbf{Artifact:} Existing benchmarks only; no release named.

\textbf{Closest oracle, and whether S1 could receive it:} \texttt{kill\_rate}.
It needs a reference implementation, not a reference specification; the
shipped code is that implementation in Lab 2, so S1 could receive it.
\end{answer}

\end{document}
